\documentclass[11pt]{article}

\usepackage[T1]{fontenc}
\usepackage{lmodern}
\usepackage{microtype}
\usepackage[a4paper,margin=32mm]{geometry}
\usepackage{amsmath,amssymb,amsthm,mathtools}
\usepackage{enumitem}
\usepackage{booktabs}
\usepackage[hidelinks]{hyperref}
\usepackage[nameinlink,noabbrev]{cleveref}

\newtheorem{theorem}{Theorem}[section]
\newtheorem{proposition}[theorem]{Proposition}
\newtheorem{lemma}[theorem]{Lemma}
\newtheorem{corollary}[theorem]{Corollary}
\theoremstyle{definition}
\newtheorem{definition}[theorem]{Definition}
\theoremstyle{remark}
\newtheorem{remark}[theorem]{Remark}

\newcommand{\T}{\mathbb{R}/\mathbb{Z}}
\newcommand{\e}{\mathrm{e}}
\newcommand{\E}{\mathbb{E}}
\newcommand{\Prob}{\mathbb{P}}
\newcommand{\RePart}{\operatorname{Re}}
\newcommand{\Hom}{\operatorname{Hom}}
\newcommand{\dist}[1]{\left\lVert #1\right\rVert}
\newcommand{\prodprime}{\mathop{\prod}\limits^{\prime}}

\title{Reciprocal Sums of Distinct Squarefree Semiprimes}
\author{Yuren Tang}
\date{}

\begin{document}

\maketitle

\begin{abstract}
We prove that a positive rational number is a finite sum of reciprocals of distinct integers, each the product of two distinct primes, if and only if its reduced denominator is squarefree.  For the converse, we first represent $1/b$ for each squarefree $b$ by a finite Fourier selection argument on a complete family of semiprime denominators.  One block of large primes serves as an anchor: low-energy anchor assignments synchronize to a single integer frequency, while the remaining prime coordinates concentrate in rows around canonical decoders.  An exact weighted product-fibre inequality compresses the Fourier sum to a decoded skeleton without discarding the lower--lower factors needed to damp intermediate frequencies.  A positive Gaussian major contribution then dominates four disjoint minor ranges.  Character orthogonality gives the desired subset sum modulo one, and a strict no-wrap bound converts it to an exact rational identity.  Avoidance and induction on the numerator complete the characterization.
\end{abstract}

\section{Introduction}
\label{sec:introduction}

An integer is a \emph{squarefree semiprime} if it is the product of two distinct primes.  We resolve the reciprocal-representation question posed by Erd\H{o}s and Graham~\cite{erdos-graham-1980} and now recorded as Erd\H{o}s Problem~306~\cite{erdos-problem-306}.

\begin{theorem}[Squarefree-denominator characterization with finite avoidance]
\label{thm:main}
Let $q>0$ be rational.  The denominator of $q$ in lowest terms is squarefree if and only if, for every finite set $\mathcal T$ of positive integers, there is a finite set $\mathcal N$ of pairwise distinct squarefree semiprimes, disjoint from $\mathcal T$, such that
\[
 q=\sum_{n\in\mathcal N}\frac1n.
\]
\end{theorem}

Thus the visible arithmetic boundary is the simple squarefree-denominator condition, while the construction is stronger than bare existence.  It survives every finite deletion.  Partial representations may be completed, finite decompositions may be realized on disjoint supports, and finite families of representations have common refinements.  These refinements may be iterated through a countably infinite proper chain whose successive supports are pairwise disjoint.  The converse direction in \cref{thm:main} is obtained directly from the fixed-target theorem by prime dilution, without numerator induction or separate treatment of reduced denominators $1$ and $2$.

Fix a reduced rational
\[
 t=\frac ab\in(0,1),
\]
with squarefree $b$, and fix $\gamma>1$.  At scale $X$ we take primes in $[X,X^\gamma)$, designate the upper fixed-ratio block as an anchor, and use all unordered distinct prime pairs together with target rows indexed by the prime-divisor set $S_b$.  The pair load tends to
\[
 \lambda_\gamma=\frac{(\log\gamma)^2}{2}.
\]
The sharp fixed-parameter region is
\begin{equation}
\label{eq:intro-admissible-region}
 t<\lambda_\gamma<1,
 \qquad\text{equivalently}\qquad
 \exp(\sqrt{2t})<\gamma<\exp(\sqrt2).
\end{equation}
Exact centring chooses $\theta\Lambda=t$, and hence
\[
 \theta\longrightarrow
 \alpha_{t,\gamma}
 =\frac{t}{\lambda_\gamma}
 =\frac{2t}{(\log\gamma)^2}.
\]

\begin{theorem}[Parameterized direct fixed-target theorem]
\label{thm:parameterized-direct-target-intro}
Let $t=a/b\in(0,1)$ be reduced with squarefree $b$, let $\gamma$ satisfy \cref{eq:intro-admissible-region}, and let $\mathcal T$ be finite.  For all sufficiently large $X$, the parameterized family of \cref{sec:denominator-system} contains a subset $\mathcal A$, disjoint from $\mathcal T$, such that
\[
 \sum_{e\in\mathcal A}\frac1e=t.
\]
Its exact target coefficient satisfies the sharp Gaussian asymptotic of \cref{thm:sharp-target-coefficient}; no moving-target uniformity is asserted.
\end{theorem}

Put $Z=X^\gamma$ and
\[
 \tau(b)=\sum_{r\in S_b}\frac1{r^2}.
\]
The actual-family variance has the single asymptotic
\[
 \sigma_{\mathcal E}^2
 \sim\alpha_{t,\gamma}(1-\alpha_{t,\gamma})
 \left(
  \frac1{2X^2\log^2X}
  +\frac{\tau(b)}{Z\log Z}
 \right).
\]
There remains one genuine phase boundary: prime pairs supply the leading variance when $\gamma>2$, whereas the target rows supply it at and below $\gamma=2$.  This change of provider does not split the proof architecture.  A universal full-coordinate cutoff
\[
 T_0=\kappa_0\min(X^2,Z)
\]
and a prime-only Gaussian transition from $T_0$ to $X^2/4$ give one pairwise-disjoint and exhaustive six-sector partition for every fixed $\gamma>1$.

The retained coefficient yields entropy rate
\[
 H(\alpha_{t,\gamma})
 =-\alpha_{t,\gamma}\log\alpha_{t,\gamma}
  -(1-\alpha_{t,\gamma})\log(1-\alpha_{t,\gamma}).
\]
For $0<t<1/2$, the maximum $\log2$ is attained at $\gamma=\exp(2\sqrt t)$.  For $t=1/2$, $\log2$ is only the supremum as $\gamma$ approaches the forbidden upper boundary.  For $1/2<t<1$, the supremum is $H(t)$ and is likewise not attained.  The direct denominator-height form, exact-cardinality extraction, Hamming-separated families, balanced signed relations, globally proper refinement branching and transfer to arbitrary admissible rationals all inherit the parameterized theorem.

\subsection{Context}

Egyptian-fraction representations and their restricted-denominator variants have a long history; see Bloom and Elsholtz~\cite{bloom-elsholtz-2022} for a recent survey.  Barbeau gave the first known representation of $1$ by reciprocals of distinct squarefree semiprimes, using $101$ terms; Watanabe later found the currently shortest known examples, with $47$ terms~\cite{barbeau-1977,watanabe-2020}.  Butler, Erd\H{o}s and Graham proved that every positive integer is a sum of distinct unit fractions whose denominators have three distinct prime factors, and conjectured the corresponding statement with two prime factors~\cite{butler-erdos-graham-2015}.

A recent preprint of Li proves the two-prime conjecture for every positive integer.  For a general squarefree denominator $b$, it also represents $a/b$ above an explicit threshold and leaves the remaining below-threshold $\omega=2$ regime open~\cite{li-2026-semiprime}.  The theorem above settles that remaining rational case.  Li's proof is driven by additive-combinatorial subset-sum input and induction; the present argument instead constructs a fixed target directly as a positive finite Fourier coefficient and obtains the full characterization by repeating a suitably diluted target on disjoint supports.

Quantitative multiplicity is known in the unrestricted setting on its natural ambient scale.  Conlon, Fox, He, Mubayi, Pham, Suk and Verstra\"ete prove that, for every fixed positive rational $q$, there are $2^{(c_q+o(1))N}$ representations by distinct denominators at most $N$~\cite{conlon-et-al-2025}.  The quantitative results here concern the much thinner squarefree-semiprime denominator system and are measured instead by the size of the constructed restricted family.

The theorem was first obtained in the author's archived Lean~4 formalization, released as version~0.0.3~\cite{tang-erdos-306-formal}.  The present article grew out of that formal work and extracts from the same Fourier-analytic line a conceptually simpler one-anchor proof intended for ordinary mathematical reading.  The archived code does not formalize this exposition line by line, because the proof underwent substantial conceptual refinement after the first formal implementation.  The archived formalization obtains its prime-distribution input from explicit Rosser--Schoenfeld estimates, whereas the present exposition derives the corresponding prime-counting and reciprocal-mass consequences from the prime number theorem.

\subsection{External input and proof architecture}

The prime number theorem is the sole non-elementary number-theoretic input; a standard reference is Montgomery and Vaughan~\cite[Chapter~1]{montgomery-vaughan-2007}.  It supplies prime counts in fixed-ratio intervals and, through Abel summation, the reciprocal mass and square load of the prime blocks.  All problem-specific estimates---anchor synchronization, multiplicity-sensitive cyclic energy, row separation, fibre compression, target observability, the universal six-sector Fourier budget, and the no-wrap passage to exact equality---are proved here.  Standard finite-group Fourier inversion, the Chinese remainder theorem, Taylor expansion, Hoeffding's inequality, and elementary combinatorial estimates are used in their usual forms.

Finite Fourier inversion identifies the coefficient that must be made positive, but it leaves the many local CRT coordinates coupled only implicitly.  The one-anchor construction makes that coupling rigid enough to estimate.

\subsection{The parameterized one-anchor construction}

Write $Z=X^\gamma$ and put
\[
 P=\{p:X\leq p<Z\},
 \qquad
 B=\{q:Z/2\leq q<Z\}.
\]
The denominator family consists of all products $pq$ with distinct $p,q\in P$, together with $rq$ for $r\in S_b$ and $q\in B$.  Its reciprocal load $\Lambda$ tends to $\lambda_\gamma$.  Under \cref{eq:intro-admissible-region}, independent Bernoulli selectors with parameter
\[
 \theta=\frac{t}{\Lambda}
\]
have expected reciprocal sum exactly $t$ and remain in a fixed compact subinterval of $(0,1)$.

Let
\[
 L=b\prod_{p\in P}p.
\]
Character orthogonality modulo $L$ expresses the probability of selecting a subset with reciprocal sum congruent to $t$ modulo one as a finite Fourier sum whose target character is $\exp(-2\pi i a h/b)$.  The numerator $a$ enters only there.  Target-coordinate rows depend solely on the reduced denominator $b$.

The anchor block synchronizes low-energy local residues to one exact integer label $m$.  Once an anchor assignment is fixed, each remaining prime coordinate is coupled to it through a row of factors indexed by $q\in B$.  The restored multiplicity-sensitive cyclic-energy lemma supplies the lower-prime row distance, while the finitely many rows indexed by $S_b$ have a stronger direct distance.  These estimates concentrate each row near a decoder and detect target directions which the complete prime-pair family alone does not observe.

The row estimates cannot simply be multiplied by the number of anchor assignments.  The decisive deterministic step is instead the exact weighted product-fibre identity of \cref{thm:weighted-compression}.  Every factor not assigned to a row remains on the decoded skeleton.  In particular, the lower--lower prime-pair factors survive the compression and provide both the universal transition damping and the adaptive damping farther out.

The skeleton has six sectors for every fixed $\gamma>1$: the Taylor major range, the total-variance Gaussian tail, the prime-only Gaussian transition, the adaptive retained-pair range, the large coherent-label range, and the noncoherent/nondecoder remainder.  The universal cutoff $T_0$ makes these ranges disjoint and exhaustive without regime-dependent tables.

The positive coefficient gives a subset whose reciprocal sum is congruent to $t$ modulo one.  Exact equality is obtained only afterwards: every subset sum lies in $[0,\Lambda]$ with $\Lambda<1$, and $t\in(0,1)$, so no nonzero integer alias is possible.

\subsection{Terminology and conventions}

The following terms describe fixed parts of the decomposition.
\begin{description}[style=nextline,widest={anchor assignment},leftmargin=*]
\item[anchor block] the upper prime block $B$;
\item[anchor assignment] the tuple of CRT coordinates indexed by $B$;
\item[coherent label] the integer $m$ whose residues agree with every coordinate of a low-energy anchor assignment;
\item[anchor-cross row] the factors joining one non-anchor prime coordinate to all primes in $B$;
\item[decoder] a row residue minimizing the associated quadratic phase energy;
\item[fibre] the set of all row-coordinate tuples above a fixed anchor assignment;
\item[decoded skeleton] the decoder tuple retained from each fibre, together with every factor not assigned to a row;
\item[target coordinate] a CRT coordinate indexed by a member of $S_b$;
\item[transition sector] the prime-only Gaussian range $T_0<|m|\leq X^2/4$.
\end{description}

Throughout the direct analytic construction, $t$, $\gamma$, $b$ and the finite forbidden set are fixed while $X\to\infty$.  The Bernoulli modulus constant $\kappa_{\rm mod}=\kappa_{\rm mod}(t,\gamma)$, the row-tail constant $c_{\rm row}=c_{\rm row}(t,\gamma,b)$ and the decoder cutoff $\kappa_0=\kappa_0(t,\gamma,b)$ have distinct roles.  Other unnamed constants may depend on the fixed tuple $(t,\gamma,b)$ and fixed numerical choices, but none depends on $X$ or on the later Gaussian cutoff $C$.  Whenever $C$ occurs, $X\to\infty$ is taken with $C$ fixed, and only afterwards may $C\to\infty$.  No statement is uniform in a moving target or moving exponent.

\subsection{Organization}

\Cref{sec:fourier-selection} records the parameterized finite Fourier identity and separates quotient realization from exact equality.  \Cref{sec:structural-tools} proves the deterministic collision, compression, syndrome and skeleton principles.  \Cref{sec:denominator-system,sec:anchor-block,sec:fibre-decoding} construct the denominator family and establish its load, total variance, anchor rigidity, cyclic row energy, row separation and decoder identification.  \Cref{sec:decoded-skeleton,sec:major-budget} give the universal six-sector analysis and prove positivity.  \Cref{sec:exact-completion} converts that positivity into exact representation, proves the characterization by prime dilution, and develops local replacement, common refinement and countable proper refinement.  \Cref{sec:quantitative-multiplicity} proves the sharp coefficient, explicit variance corollaries, entropy optimization and quantitative consequences.  \Cref{sec:scope} records the scope and effectivity boundary.

\section{Finite Fourier selection and exactness}
\label{sec:fourier-selection}

The construction begins with a finite-group coefficient.  We record the selection identity in a form suited to a fixed rational target and keep the later passage from a quotient congruence to exact equality logically separate.

\subsection{Character orthogonality}

Let $A$ be a finite abelian group and let $\widehat A$ be its character group.  For $a,t\in A$,
\[
 \frac1{|A|}\sum_{\chi\in\widehat A}\overline{\chi(t)}\chi(a)
 =\mathbf 1_{a=t}.
\]
Consequently, if $T$ is an $A$-valued random variable, then
\begin{equation}
\label{eq:finite-fourier-coefficient}
 \Prob(T=t)
 =\frac1{|A|}\sum_{\chi\in\widehat A}
   \overline{\chi(t)}\E\chi(T).
\end{equation}
The numerator is the real nonnegative number $|A|\Prob(T=t)$, so it is enough to prove that its real part is positive.

\begin{lemma}[Independent product formula]
\label{lem:independent-product}
Let $v_j\in A$, and let $\xi_j$ be independent Bernoulli random variables with parameters $\theta_j\in[0,1]$.  For $T=\sum_j\xi_jv_j$ and every $\chi\in\widehat A$,
\[
 \E\chi(T)=\prod_j\bigl((1-\theta_j)+\theta_j\chi(v_j)\bigr).
\]
\end{lemma}

\begin{proof}
Since $\chi$ is a homomorphism, $\chi(T)=\prod_j\chi(v_j)^{\xi_j}$.  Independence gives the product of expectations.
\end{proof}

\subsection{Reciprocal subset sums at a fixed target}

Let $\mathcal E$ be a finite set of positive integers and let $L$ be divisible by every $e\in\mathcal E$.  Associate to $e$ the residue $v_e=L/e\pmod L$.  Let
\[
 t=\frac ab\in(0,1)
\]
be reduced, with $b\mid L$, and let $\xi_e$ be independent Bernoulli variables with common parameter $\theta$.  Then
\[
 T=\sum_{e\in\mathcal E}\xi_ev_e
\]
encodes the subset reciprocal sum modulo one: the congruence $T=aL/b\pmod L$ is equivalent to
\[
 \sum_{e\in\mathcal E}\frac{\xi_e}{e}\equiv\frac ab\pmod1.
\]
Writing characters as $h\mapsto\exp(2\pi i h\,\cdot/L)$, \cref{eq:finite-fourier-coefficient} gives
\begin{equation}
\label{eq:reciprocal-fourier}
 L\Prob\left(T=\frac{aL}{b}\right)
 =\sum_{h\bmod L}
   \exp\left(-\frac{2\pi i a h}{b}\right)
   \prod_{e\in\mathcal E}
   \left((1-\theta)+\theta\exp\left(\frac{2\pi i h}{e}\right)\right).
\end{equation}
We shall prove that the real part of the right side is positive.

The linear phase is centred by choosing
\begin{equation}
\label{eq:theta-centering}
 \theta\sum_{e\in\mathcal E}\frac1e=t.
\end{equation}
For the parameterized semiprime family this is $\theta=t/\Lambda$, where $\Lambda$ is the total reciprocal load.  Apart from its scalar effect through this centring parameter, the numerator $a$ enters the Fourier geometry only through the target character in \cref{eq:reciprocal-fourier}; the CRT rows which observe the target depend only on the prime divisors of $b$.

\subsection{Quotient realization and ambient equality}

Fourier positivity in a quotient proves only a quotient target.  We therefore record the exactness step separately.

Let $M$ be an abelian group, let $\pi:M\to A$ be a homomorphism to a finite abelian group, let $T_M$ be an $M$-valued random variable, and let $t_M\in M$.  Positivity of the coefficient of $\pi(t_M)$ gives $\Prob(\pi(T_M)=\pi(t_M))>0$.

\begin{proposition}[Exactness interfaces]
\label{prop:exactness-interfaces}
Each of the following is sufficient for $\Prob(T_M=t_M)>0$.
\begin{enumerate}[label=(\roman*)]
\item There is a set $C\subset M$ containing $t_M$ and the support of $T_M$ such that $\pi|_C$ is injective.
\item The support of $T_M$ and $t_M$ lie in one fundamental domain for $\ker\pi$.
\item One has the strict measurable alias inequality
\[
 \Prob\bigl(\pi(T_M)=\pi(t_M)\bigr)
 >\Prob\bigl(T_M\in(t_M+\ker\pi)\setminus\{t_M\}\bigr).
\]
\end{enumerate}
\end{proposition}

\begin{proof}
Under~(i) or~(ii), the only point in the relevant quotient fibre which can lie in the support is $t_M$.  For~(iii), decompose the quotient event into the target event and the disjoint nonzero-alias event.
\end{proof}

In the present construction, all reciprocal subset sums and the target lie in $M=L^{-1}\mathbb Z$, and reduction modulo $\mathbb Z$ gives $A\cong\mathbb Z/L\mathbb Z$.  We use only the fundamental-domain criterion.  The total load satisfies $\Lambda<1$, so every subset sum lies in $[0,1)$, as does the fixed target $t$.  Hence congruence modulo one is exact equality.  This no-wrap argument is invoked only after positivity has been proved.

\subsection{Chinese-remainder coordinates}

The period $L$ is squarefree.  Hence the Chinese remainder theorem identifies
\[
 \mathbb Z/L\mathbb Z\cong\prod_{p\mid L}\mathbb Z/p\mathbb Z.
\]
A chosen block of large prime coordinates forms the anchor block.  The remaining prime coordinates are grouped into anchor-cross rows whose kernels depend on the anchor assignment.  Denominator factors involving two anchor primes contribute to the anchor weight; factors involving one anchor prime and one row prime contribute to that row; every other factor remains in a residual term on the decoded skeleton.  This is an exact partition of the original Fourier product.

\section{Structural tools}
\label{sec:structural-tools}

This section isolates the deterministic statements used by the arithmetic construction.  The complete-family estimate controls reciprocal capacity; the product-fibre theorem performs the compression to one decoder tuple per anchor; the syndrome theorem detects hidden finite-group directions; and the final proposition assembles a positive skeleton contribution against finitely many error lanes.  The prime-specific rigidity, row distance and minor-frequency estimates are supplied only in the later sections.

\subsection{Complete weighted families}

Let $I$ be a finite set and let $u_i\geq0$ for $i\in I$.  Put
\[
 U=\sum_{i\in I}u_i,
 \qquad
 V=\sum_{i\in I}u_i^2.
\]
For $k\geq0$, write
\[
 e_k(u)=\sum_{\substack{J\subset I\\ |J|=k}}\prod_{j\in J}u_j,
\]
with $e_0(u)=1$ and $e_k(u)=0$ for $k>|I|$.

\begin{theorem}[Complete-family collision bound]
\label{thm:complete-family-collision}
For every integer $k\geq2$,
\[
 0\leq \frac{U^k}{k!}-e_k(u)
 \leq \frac{U^{k-2}V}{2(k-2)!}.
\]
\end{theorem}

\begin{proof}
Expand
\[
 U^k=\sum_{(i_1,\ldots,i_k)\in I^k}\prod_{a=1}^k u_{i_a}.
\]
The ordered tuples with pairwise distinct coordinates contribute exactly $k!e_k(u)$.  Thus $U^k-k!e_k(u)$ is the nonnegative total weight of tuples containing a collision.

For fixed positions $1\leq a<b\leq k$, the total weight of tuples satisfying $i_a=i_b$ is $VU^{k-2}$.  The union bound over the $\binom{k}{2}$ pairs of positions gives
\[
 U^k-k!e_k(u)
 \leq \binom{k}{2}VU^{k-2}.
\]
Division by $k!$ proves the assertion.
\end{proof}

\begin{corollary}[Relative and asymptotic forms]
\label{cor:complete-family-relative}
If $U>0$, then
\[
 1-\binom{k}{2}\frac{V}{U^2}
 \leq \frac{k!e_k(u)}{U^k}
 \leq1.
\]
Consequently, for a sequence of weighted families:
\begin{enumerate}[label=(\roman*)]
\item for fixed $k$, the condition $V/U^2\to0$ implies
\[
 e_k(u)=\frac{U^k}{k!}(1+o(1));
\]
\item for moving integers $k=k_n$, the same conclusion holds whenever
\[
 k^2\frac{V}{U^2}\longrightarrow0.
\]
\end{enumerate}
If $U=0$, every $u_i$ vanishes and the assertions for $k\geq1$ are immediate.
\end{corollary}

For pairs there is an exact identity rather than merely an estimate.

\begin{corollary}[Pair identity]
\label{cor:pair-identity}
One has
\[
 e_2(u)=\frac{U^2-V}{2}.
\]
\end{corollary}

\begin{proof}
This is the identity
\[
 U^2=\sum_i u_i^2+2\sum_{i<j}u_iu_j.
\]
\end{proof}

Only \cref{cor:pair-identity} is needed for the dense semiprime family.  The higher-uniformity statement records the natural collision-free range detected by the first two power sums; it does not provide the rigidity or decoding theory required by a higher-uniformity application.

\subsection{Exact weighted product-fibre compression}

Let $Y$ be finite and let $R$ be a finite, possibly empty, coordinate set.  For each $r\in R$, let $X_r$ be finite and nonempty, and put
\[
 X=\prod_{r\in R}X_r.
\]
For each $y\in Y$, suppose that
\begin{itemize}
\item $W_y\geq0$;
\item $a_{r,y}:X_r\to[0,\infty)$ for $r\in R$;
\item a decoder $d_r(y)\in X_r$ has been chosen;
\item $H_y:X\to\mathbb C$ satisfies $|H_y(x)|\leq1$.
\end{itemize}
Define
\[
 F(y,x)=W_y\prod_{r\in R}a_{r,y}(x_r)H_y(x),
\]
and
\[
 \alpha_{r,y}=a_{r,y}(d_r(y)),
 \qquad
 \beta_{r,y}=\sum_{x_r\neq d_r(y)}a_{r,y}(x_r),
 \qquad
 d(y)=(d_r(y))_{r\in R}.
\]
Empty products are $1$.

\begin{theorem}[Exact weighted product-fibre compression]
\label{thm:weighted-compression}
There is a complex number $E_{\mathrm{fib}}$ such that
\begin{equation}
\label{eq:fibre-decomposition}
 \sum_{y\in Y}\sum_{x\in X}F(y,x)
 =\sum_{y\in Y}F(y,d(y))+E_{\mathrm{fib}},
\end{equation}
and
\begin{equation}
\label{eq:exact-fibre-error}
 |E_{\mathrm{fib}}|
 \leq
 \sum_{y\in Y}W_y
 \left\{
   \prod_{r\in R}(\alpha_{r,y}+\beta_{r,y})
   -\prod_{r\in R}\alpha_{r,y}
 \right\}.
\end{equation}
The bound is sharp among deterministic bounds depending only on the displayed values $W_y$, $\alpha_{r,y}$ and $\beta_{r,y}$.
\end{theorem}

\begin{proof}
Fix $y$.  Separate the decoder tuple $d(y)$ from the other tuples.  Since $|H_y(x)|\leq1$,
\[
 \left|\sum_{x\neq d(y)}F(y,x)\right|
 \leq W_y\sum_{x\neq d(y)}\prod_r a_{r,y}(x_r).
\]
The full nonnegative product sum factorizes as
\[
 \sum_{x\in X}\prod_r a_{r,y}(x_r)
 =\prod_r\sum_{x_r\in X_r}a_{r,y}(x_r)
 =\prod_r(\alpha_{r,y}+\beta_{r,y}),
\]
whereas the decoder tuple has product weight $\prod_r\alpha_{r,y}$.  Their difference is therefore the exact off-decoder product mass.  Summing over $y$ proves \cref{eq:fibre-decomposition,eq:exact-fibre-error}.

For sharpness, choose the off-decoder values of $H_y$ to have a common unit phase.  The complex remainder then attains the entire off-decoder product mass.
\end{proof}

Several boundary cases are included in the statement.  A zero anchor weight contributes nothing.  If some $\alpha_{r,y}$ vanishes, \cref{thm:weighted-compression} remains valid but no division by that decoder weight is permitted.  If $R$ is empty, then $X$ is a singleton and $E_{\mathrm{fib}}=0$.  Decoder uniqueness is not assumed: other minimizers remain in $\beta_{r,y}$ and are charged as off-decoder mass.  Complex row kernels are handled by placing their magnitudes in $a_{r,y}$ and retaining every phase, together with every factor not assigned to a row, in $H_y$.

When all decoder weights are positive, a normalized form is available.

\begin{corollary}[Anchor-dependent normalized average]
\label{cor:normalized-fibre}
Assume $\alpha_{r,y}>0$ for all $r,y$, and put
\[
 \varepsilon_{r,y}=\frac{\beta_{r,y}}{\alpha_{r,y}},
 \qquad
 A_y=W_y\prod_r\alpha_{r,y},
 \qquad
 \Delta_y=\sum_r\varepsilon_{r,y}.
\]
Then
\[
 |E_{\mathrm{fib}}|
 \leq\sum_yA_y\left\{\prod_r(1+\varepsilon_{r,y})-1\right\}
 \leq\sum_yA_y(\e^{\Delta_y}-1).
\]
\end{corollary}

This is an anchor-dependent weighted average.  No sequential conditional kernel is asserted.  In the semiprime application we use the unnormalized theorem, so zero decoder weights never create a normalization issue.

\subsection{Syndrome separation for hidden finite groups}

Let $H$ be a finite abelian group, let $J$ be a finite sensor set, and let
\[
 \psi_j\in\Hom(H,\T)
 \qquad (j\in J).
\]
For $t\in\T$, let $\dist{t}$ be its distance to $0$.  When $H$ is nontrivial, define
\[
 \delta^2=\min_{0\neq h\in H}\sum_{j\in J}\dist{\psi_j(h)}^2.
\]
For shifts $a_j\in\T$, put
\[
 E_a(h)=\sum_{j\in J}\dist{a_j+\psi_j(h)}^2.
\]

\begin{theorem}[Nearest-syndrome separation]
\label{thm:syndrome-separation}
Assume $H$ is nontrivial and choose a minimizer $h_a$ of $E_a$.  Then for every $h\neq h_a$,
\[
 E_a(h)\geq\frac{\delta^2}{4}.
\]
\end{theorem}

\begin{proof}
For every $j$, the circle metric and the homomorphism property give
\[
 \dist{\psi_j(h-h_a)}
 \leq \dist{a_j+\psi_j(h)}+\dist{a_j+\psi_j(h_a)}.
\]
After squaring and applying $(u+v)^2\leq2u^2+2v^2$, we obtain
\[
 \delta^2
 \leq\sum_j\dist{\psi_j(h-h_a)}^2
 \leq2E_a(h)+2E_a(h_a).
\]
Since $h_a$ minimizes $E_a$, one has $E_a(h_a)\leq E_a(h)$, and hence $\delta^2\leq4E_a(h)$.
\end{proof}

\begin{corollary}[Gaussian-modulus decoder tail]
\label{cor:syndrome-tail}
Suppose $K_j:\T\to\mathbb C$ satisfy
\[
 |K_j(t)|\leq\exp(-\kappa\dist{t}^2)
\]
for some $\kappa>0$.  Then
\[
 \sum_{h\neq h_a}\prod_{j\in J}|K_j(a_j+\psi_j(h))|
 \leq(|H|-1)\exp(-\kappa\delta^2/4).
\]
\end{corollary}

\begin{proposition}[Qualitative observability]
\label{prop:qualitative-observability}
For nontrivial finite $H$, the following are equivalent:
\begin{enumerate}[label=(\roman*)]
\item $\delta>0$;
\item the syndrome map
\[
 \Sigma:H\longrightarrow\T^J,
 \qquad
 \Sigma(h)=(\psi_j(h))_{j\in J},
\]
is injective;
\item $\bigcap_{j\in J}\ker\psi_j=\{0\}$.
\end{enumerate}
\end{proposition}

\begin{proof}
Injectivity of $\Sigma$ is equivalent to triviality of its kernel, which is the common kernel of the sensors.  Since $H$ is finite, injectivity makes every nonzero syndrome have positive squared norm and hence gives $\delta>0$.  Conversely, $\delta=0$ means that some nonzero element lies in every sensor kernel.
\end{proof}

For one finite system, $\delta>0$ is a qualitative injectivity statement.  For a sequence of systems, an explicit scale-dependent lower bound for $\delta$ is the quantitative input.  The theorem gives a sufficient absolute-value decoding estimate; it does not classify arguments which exploit cancellation inside fibres.  It also does not identify a prescribed decoder without an application-specific comparison of the candidate energy with the competing syndrome distance.  If $H$ is trivial, there is no hidden-coordinate obligation and no minimum over nonzero elements is defined.

We use the following coverage principle.  In each frequency lane, every hidden quotient direction must be fixed by the anchor, resolved by a decoded fibre, or detected by a retained residual factor which damps that lane.  Thus observability is a distinct obligation, although in a concrete lane it may be supplied by row separation, decoder identification or residual damping.

\subsection{Decoded-skeleton positivity}

Retain the notation of \cref{thm:weighted-compression}.  Partition the anchor set into one major set and finitely many minor lanes,
\[
 Y=Y_{\mathrm{maj}}\sqcup Y_1\sqcup\cdots\sqcup Y_s.
\]
No abstract distinction between ``middle'', ``outer'' or ``energetic'' lanes is imposed.

\begin{proposition}[Decoded-skeleton positivity]
\label{prop:decoded-skeleton-positivity}
Put
\[
 E_{\mathrm{prod}}=
 \sum_yW_y\left\{
   \prod_r(\alpha_{r,y}+\beta_{r,y})
   -\prod_r\alpha_{r,y}
 \right\}.
\]
Then
\begin{align*}
 \RePart\sum_{y,x}F(y,x)
 \geq{}&
 \RePart\sum_{y\in Y_{\mathrm{maj}}}F(y,d(y))\\
 &-\sum_{j=1}^s\left|\sum_{y\in Y_j}F(y,d(y))\right|
 -E_{\mathrm{prod}}.
\end{align*}
In particular, the full sum has positive real part if the displayed major contribution strictly exceeds the sum of the minor-lane bounds and $E_{\mathrm{prod}}$.
\end{proposition}

\begin{proof}
By \cref{thm:weighted-compression},
\[
 \sum_{y,x}F(y,x)=\sum_yF(y,d(y))+E_{\mathrm{fib}},
 \qquad
 |E_{\mathrm{fib}}|\leq E_{\mathrm{prod}}.
\]
Split the skeleton sum according to the displayed partition, take the real part of the major lane, and bound every minor-lane sum and the fibre remainder in absolute value.
\end{proof}

For a sequence of systems, let $M_n>0$ eventually.  If the major real part is at least $cM_n$ for some fixed $c>0$, while the total of all minor-lane and product-fibre errors is $o(M_n)$, then the full sum has positive real part for all sufficiently large $n$.

A coherent label map is optional in \cref{prop:decoded-skeleton-positivity}; it is introduced only when it improves the estimates on the skeleton lanes.  The obligations used to prove the lane bounds remain logically distinct, but one application theorem may discharge several of them on disjoint lanes.  In particular, the same anchor theorem may provide both a weighted partition and a large-label tail without double-counting.

The proposition proves positivity of the finite Fourier sum to which it is applied.  If this sum is the numerator of a finite-group coefficient, it gives a realized quotient target.  Exact ambient equality is obtained afterwards through one of the interfaces in \cref{prop:exactness-interfaces}.  For the semiprime theorem we use only deterministic no-wrap.

\section{The parameterized semiprime denominator system}
\label{sec:denominator-system}

Fix throughout a reduced target
\[
 t=\frac ab\in(0,1)
\]
with squarefree $b$, a fixed anchor ratio $\eta\in(0,1)$, a finite forbidden set $\mathcal T$, and a fixed exponent $\gamma>1$.  Reducedness and $0<t<1$ imply $b>1$.  Thus the set $S_b$ of prime divisors of $b$ is nonempty and
\begin{equation}
\label{eq:tau-positive}
 \tau(b):=\sum_{r\in S_b}\frac1{r^2}>0.
\end{equation}
Put
\begin{equation}
\label{eq:lambda-gamma}
 \lambda_\gamma=\frac{(\log\gamma)^2}{2}.
\end{equation}
We assume the exact inequalities for the centred no-wrap construction
\begin{equation}
\label{eq:admissible-region}
 t<\lambda_\gamma<1.
\end{equation}
Equivalently,
\begin{equation}
\label{eq:admissible-gamma}
 \exp(\sqrt{2t})<\gamma<\exp(\sqrt2).
\end{equation}
All asymptotic statements are for $X\to\infty$ with $t$, $\gamma$, $\eta$, $b$ and $\mathcal T$ fixed.  Unless stated otherwise, implied constants may depend on the fixed tuple $(t,\gamma,\eta,b)$ and on fixed numerical choices, but not on $X$ or on the later Gaussian cutoff $C$.  The estimates below are uniform when $\eta$ ranges in a fixed compact subset of $(0,1)$.

\subsection{Prime blocks and their power sums}

Put
\begin{equation}
\label{eq:prime-blocks}
 Z=X^\gamma,
 \qquad
 P=\{p\text{ prime}:X\leq p<Z\},
 \qquad
 B=\{q\text{ prime}:\eta Z\leq q<Z\},
\end{equation}
and let $S_b$ be the set of prime divisors of $b$.  We take $X$ so large that every prime in $P$ exceeds every prime divisor of $b$, and every denominator constructed below exceeds every element of $\mathcal T$.

We use the prime number theorem in the forms
\[
 \pi(y)\sim\frac{y}{\log y},
 \qquad
 \vartheta(y):=\sum_{p\leq y}\log p=y+o(y),
\]
together with their standard partial-summation consequences; see Montgomery and Vaughan~\cite[Chapter~1]{montgomery-vaughan-2007}.

\begin{lemma}[Prime supply in fixed-ratio intervals]
\label{lem:fixed-ratio-prime-supply}
For fixed $0<u<v$, as $y\to\infty$,
\[
 \#\{p\text{ prime}:uy\leq p<vy\}
 =\bigl(v-u+o_{u,v}(1)\bigr)\frac{y}{\log y}.
\]
The error is uniform when $(u,v)$ ranges over a compact subset of $\{(u,v):0<u<v\}$.
\end{lemma}

\begin{proof}
Apply the prime number theorem at $uy$ and $vy$.  Replacing $\log(uy)$ and $\log(vy)$ by $\log y$ introduces only a relative $o_{u,v}(1)$, uniformly on compact parameter sets.
\end{proof}

Define
\[
 S_X=\sum_{p\in P}\frac1p,
 \qquad
 U_X=\sum_{p\in P}\frac1{p^2},
 \qquad
 M_B=|B|,
\]
and
\[
 V_B=\sum_{q\in B}\frac1q,
 \qquad
 W_B=\sum_{q\in B}\frac1{q^2}.
\]

\begin{lemma}[Parameterized prime reciprocal estimates]
\label{lem:reciprocal-prime-estimates}
As $X\to\infty$,
\begin{align}
 S_X&=\log\gamma+o(1),
 \label{eq:SX-asymptotic}\\
 U_X&\sim\frac1{X\log X},
 \label{eq:UX-asymptotic}\\
 |P|&\sim\frac{Z}{\log Z},
 \label{eq:P-cardinality}\\
 M_B&=\bigl(1-\eta+o_\eta(1)\bigr)\frac{Z}{\log Z},
 \label{eq:B-cardinality}\\
 V_B&=\frac{\log(1/\eta)+o_\eta(1)}{\log Z},
 \label{eq:B-reciprocal-mass}\\
 W_B&=\frac{\eta^{-1}-1+o_\eta(1)}{Z\log Z}.
 \label{eq:B-power-sums}
\end{align}
\end{lemma}

\begin{proof}
Abel summation gives
\[
 \sum_{X\leq p<X^\gamma}\frac1p
 =\int_X^{X^\gamma}\frac{du}{u\log u}+o(1)
 =\log\gamma+o(1).
\]
A second partial summation gives
\[
 \sum_{X\leq p<X^\gamma}\frac1{p^2}
 =\int_X^{X^\gamma}\frac{du}{u^2\log u}(1+o(1))
 \sim\frac1{X\log X};
\]
the upper endpoint is negligible because $\gamma>1$ is fixed.  The cardinality statements follow from \cref{lem:fixed-ratio-prime-supply}.  Partial summation on $[\eta Z,Z)$ gives
\[
 \sum_{\eta Z\leq q<Z}\frac1q
 =\int_{\eta Z}^{Z}\frac{du}{u\log u}+o_\eta(1/\log Z)
 =\frac{\log(1/\eta)+o_\eta(1)}{\log Z}
\]
and
\[
 \sum_{\eta Z\leq q<Z}\frac1{q^2}
 =\int_{\eta Z}^{Z}\frac{du}{u^2\log u}(1+o_\eta(1))
 =\frac{\eta^{-1}-1+o_\eta(1)}{Z\log Z}.
\]
All fixed-ratio estimates are uniform for $\eta$ in compact subsets of $(0,1)$.
\end{proof}

\subsection{The complete pair family and target rows}

Define
\[
 \mathcal E_{\mathrm{pair}}
 =\{pq:p,q\in P,\ p<q\},
 \qquad
 \mathcal E_b
 =\{rq:r\in S_b,\ q\in B\},
\]
and
\begin{equation}
\label{eq:E-and-L}
 \mathcal E=\mathcal E_{\mathrm{pair}}\sqcup\mathcal E_b,
 \qquad
 L=b\prod_{p\in P}p.
\end{equation}
The disjoint union follows from unique factorization.  Put
\[
 \Lambda=\sum_{e\in\mathcal E}\frac1e.
\]

\begin{proposition}[Parameterized arithmetic capacity]
\label{prop:arithmetic-capacity}
For all sufficiently large $X$:
\begin{enumerate}[label=(\roman*)]
\item every $e\in\mathcal E$ is a squarefree semiprime;
\item the elements of $\mathcal E$ are pairwise distinct, avoid $\mathcal T$, and divide $L$;
\item one has
\begin{equation}
\label{eq:load-limit}
 \Lambda\longrightarrow\lambda_\gamma;
\end{equation}
more precisely, for fixed constants
\[
 t<\lambda_-<\lambda_\gamma<\lambda_+<1,
\]
one has $\lambda_-<\Lambda<\lambda_+$ for all sufficiently large $X$;
\item the complete square load satisfies
\begin{equation}
\label{eq:complete-square-load}
 \sum_{e\in\mathcal E}\frac1{e^2}
 =\frac{1+o(1)}{2X^2\log^2X}
  +\frac{(\eta^{-1}-1)\tau(b)+o_{b,\eta}(1)}{Z\log Z}.
\end{equation}
\end{enumerate}
\end{proposition}

\begin{proof}
Only the load and square load require proof.  The exact pair identity \cref{cor:pair-identity} gives
\begin{equation}
\label{eq:pair-load}
 \Lambda_{\mathrm{pair}}
 =\frac12(S_X^2-U_X)
 \longrightarrow\frac{(\log\gamma)^2}{2}=\lambda_\gamma.
\end{equation}
The target-row load is
\[
 \Lambda_b
 =\left(\sum_{r\in S_b}\frac1r\right)V_B
 =\left(\sum_{r\in S_b}\frac1r\right)
   \frac{\log(1/\eta)+o_\eta(1)}{\log Z}
 =o_{b,\eta}(1),
\]
which proves \cref{eq:load-limit} and the fixed load interval.

For the pair square load, apply the pair identity to $u_p=p^{-2}$:
\[
 \sum_{p<q}\frac1{p^2q^2}
 =\frac12\left(U_X^2-\sum_{p\in P}\frac1{p^4}\right)
 \sim\frac1{2X^2\log^2X}.
\]
Indeed the fourth-power diagonal is $O(1/(X^3\log X))$.  The target rows contribute
\[
 \sum_{r\in S_b}\sum_{q\in B}\frac1{r^2q^2}
 =\tau(b)W_B
 =\frac{(\eta^{-1}-1)\tau(b)+o_{b,\eta}(1)}{Z\log Z}.
\]
This proves \cref{eq:complete-square-load} without assuming which summand dominates.
\end{proof}

\begin{proposition}[Family size and modulus]
\label{prop:family-size-modulus}
One has
\begin{equation}
\label{eq:family-size-modulus}
 M_X:=|\mathcal E|
 =(1+o(1))\frac{Z^2}{2\log^2Z},
 \qquad
 \log L=(1+o(1))Z=o(M_X).
\end{equation}
The target rows have
\[
 |\mathcal E_b|
 =\bigl(|S_b|(1-\eta)+o_{b,\eta}(1)\bigr)\frac{Z}{\log Z}
 =O_{b,\eta}(Z/\log Z)
\]
elements and do not affect the leading family size.  Every denominator is at most $Z^2$ for large $X$.
\end{proposition}

\begin{proof}
The pair family has $\binom{|P|}{2}=(1+o(1))Z^2/(2\log^2Z)$ elements by \cref{eq:P-cardinality}, while the displayed target-row count follows from \cref{eq:B-cardinality}.  This proves the formula for $M_X$.

For the modulus, the prime number theorem in the form $\vartheta(y)=y+o(y)$ gives
\begin{align*}
 \log L
 &=\log b+\sum_{p\in P}\log p\\
 &=\vartheta(Z)-\vartheta(X)+O_b(1)\\
 &=(1+o(1))Z.
\end{align*}
The harmless $O_b(1)$ accounts for the endpoint convention and the fixed factor $b$.  Since $\gamma>1$, one has $X=o(Z)$.  Finally,
\[
 \frac{\log L}{M_X}
 =(2+o(1))\frac{\log^2Z}{Z}\longrightarrow0,
\]
so $\log L=o(M_X)$.
\end{proof}

\subsection{Exact centring and the variance boundary}

Set
\begin{equation}
\label{eq:theta-definition}
 \theta=\frac{t}{\Lambda}.
\end{equation}
By \cref{prop:arithmetic-capacity}, $\theta$ lies in a fixed compact interval $I_{t,\gamma,\eta}\subset(0,1)$ and
\begin{equation}
\label{eq:theta-limit}
 \theta\longrightarrow
 \alpha_{t,\gamma}
 :=\frac{t}{\lambda_\gamma}
 =\frac{2t}{(\log\gamma)^2}.
\end{equation}
Moreover,
\begin{equation}
\label{eq:exact-centering-load}
 \theta\sum_{e\in\mathcal E}\frac1e=t,
 \qquad
 \Lambda<1.
\end{equation}
Define the actual-family variance by
\begin{equation}
\label{eq:sigma-E}
 \sigma_{\mathcal E}^2
 =\theta(1-\theta)\sum_{e\in\mathcal E}\frac1{e^2}.
\end{equation}

\begin{proposition}[Total actual-family variance]
\label{prop:variance-regimes}
Let $\alpha=\alpha_{t,\gamma}$.  Then
\begin{equation}
\label{eq:total-variance}
 \sigma_{\mathcal E}^2
 \sim \alpha(1-\alpha)
 \left(
  \frac1{2X^2\log^2X}
  +\frac{(\eta^{-1}-1)\tau(b)}{Z\log Z}
 \right).
\end{equation}
\end{proposition}

\begin{proof}
Combine \cref{eq:theta-limit,eq:complete-square-load,eq:sigma-E}.  Both summands are retained, so the statement is valid without first deciding which contribution dominates.
\end{proof}

\begin{corollary}[Explicit variance regimes]
\label{cor:variance-regimes}
With $\alpha=\alpha_{t,\gamma}$:
\begin{enumerate}[label=(\Alph*)]
\item If $\gamma>2$, the pair square load dominates and
\begin{equation}
\label{eq:variance-gamma-large}
 \sigma_{\mathcal E}^2
 \sim\frac{\alpha(1-\alpha)}{2X^2\log^2X}.
\end{equation}
\item If $\gamma=2$, the target-row square load dominates the pair square load by a factor of order $\log X$, and
\begin{equation}
\label{eq:variance-gamma-two}
 \sigma_{\mathcal E}^2
 \sim\frac{\alpha(1-\alpha)(\eta^{-1}-1)\tau(b)}{Z\log Z}.
\end{equation}
\item If $1<\gamma<2$, the target-row square load dominates by a power of $X$, and the same asymptotic \cref{eq:variance-gamma-two} holds.
\end{enumerate}
\end{corollary}

\begin{proof}
The ratio of the target-row summand in \cref{eq:total-variance} to the pair summand is asymptotic, up to a fixed positive factor depending on $(\eta,b)$, to
\[
 \frac{X^2\log^2X}{Z\log Z}
 \asymp X^{2-\gamma}\log X.
\]
It tends to $0$ for $\gamma>2$, to infinity logarithmically for $\gamma=2$, and to infinity polynomially for $\gamma<2$.  Positivity of $(\eta^{-1}-1)\tau(b)$ follows from $\eta\in(0,1)$ and \cref{eq:tau-positive}; the case $b=1$ cannot occur in the direct setting $t\in(0,1)$.
\end{proof}

Thus $\gamma=2$ is a genuine variance boundary for every fixed $\eta\in(0,1)$, but it does not require a separate proof architecture: every later estimate is stated first with the total variance \cref{eq:total-variance}.

The finite Fourier numerator is
\begin{equation}
\label{eq:Fh-definition}
 F(h)=\exp\left(-\frac{2\pi i a h}{b}\right)
 \prod_{e\in\mathcal E}
 \left((1-\theta)+\theta\exp\left(\frac{2\pi i h}{e}\right)\right).
\end{equation}
Compactness of $I_{t,\gamma,\eta}$ gives a constant
\[
 \kappa_{\rm mod}=\kappa_{\rm mod}(t,\gamma,\eta)>0
\]
such that
\begin{equation}
\label{eq:bernoulli-modulus}
 |(1-\theta)+\theta\exp(2\pi i u)|
 \leq\exp(-\kappa_{\rm mod}\dist{u}^2).
\end{equation}
Consequently,
\begin{equation}
\label{eq:global-Fourier-energy}
 |F(h)|\leq\exp(-\kappa_{\rm mod}Q_{\mathcal E}(h)),
 \qquad
 Q_{\mathcal E}(h)=\sum_{e\in\mathcal E}\dist{h/e}^2.
\end{equation}

\subsection{Exact factor partition}

Because $b$ and $\prod_{p\in P}p$ are coprime and squarefree, the period $L$ is squarefree.  Chinese-remainder coordinates split the frequency into those indexed by $B$, by $P\setminus B$, and by $S_b$.  The denominator factors are assigned exactly once as follows:
\begin{center}
\begin{tabular}{ll}
\toprule
Denominator class & role \\
\midrule
$q q'$ with $q,q'\in B$ & internal anchor factor \\
$r q$ with $r\in P\setminus B$, $q\in B$ & lower-prime anchor-cross row $r$ \\
$r q$ with $r\in S_b$, $q\in B$ & target-coordinate row $r$ \\
$p p'$ with $p,p'\in P\setminus B$ & retained skeleton factor \\
\bottomrule
\end{tabular}
\end{center}
The target character and every phase not absorbed into a row remain in the complex residual factor.  The four denominator classes are disjoint and cover $\mathcal E$.

\input{sections/05_top_anchor}
\section{Fibre decoding and target observability}
\label{sec:fibre-decoding}

Fix an anchor assignment.  The remaining CRT coordinates are indexed by
\[
 R_{\mathrm{rows}}=(P\setminus B)\sqcup S_b.
\]
For $r\in R_{\mathrm{rows}}$, a residue $x\pmod r$, and $q\in B$, let $H_{rq}(x,a_q)$ be the centred CRT lift of $x\pmod r$ and $a_q\pmod q$, and put
\[
 u_{rq}(x)=\frac{H_{rq}(x,a_q)}{rq}\pmod1.
\]
If the row residue changes by $d$, then
\begin{equation}
\label{eq:row-shift}
 u_{rq}(x+d)-u_{rq}(x)=dq^{-1}/r\pmod1.
\end{equation}
Thus the separation of two possible row residues is independent of the anchor shift.

\subsection{Cyclic row distance}

\begin{lemma}[Multiplicity-sensitive cyclic energy]
\label{lem:cyclic-energy}
Let $c_1,\ldots,c_M$ be nonzero residues modulo a prime $r$, each occurring at most $\mu$ times.  If $M/\mu$ is sufficiently large, then for every $d\not\equiv0\pmod r$,
\[
 \sum_{j=1}^M\dist{dc_j/r}^2
 \geq c\frac{M^3}{\mu^2r^2}
\]
for an absolute constant $c>0$.
\end{lemma}

\begin{proof}
Multiplication by the nonzero residue $d$ permutes the nonzero residue classes modulo $r$ and therefore preserves the multiplicity bound.  Put
\[
 u=\left\lfloor\frac{M}{4\mu}\right\rfloor.
\]
For sufficiently large $M/\mu$, one has $u\geq M/(8\mu)\geq1$.  At most $2\mu u\leq M/2$ entries can lie in the $2u$ nonzero residue classes closest to zero on the circle.  Every remaining entry has circle distance at least
\[
 \frac{u}{r}\geq\frac{M}{8\mu r}.
\]
At least $M/2$ entries therefore contribute the square of this quantity, which proves the stated lower bound.
\end{proof}

For a row $r$, define
\[
 D_r=\min_{d\not\equiv0\,(r)}\sum_{q\in B}\dist{dq^{-1}/r}^2.
\]

\begin{proposition}[Lower-prime row distance]
\label{prop:lower-row-distance}
Uniformly for every prime $r\in P\setminus B$,
\[
 D_r\gg\frac{Z}{\log^3Z}.
\]
\end{proposition}

\begin{proof}
Inversion permutes the nonzero residue classes modulo $r$.  The interval $[Z/2,Z)$ contains at most
\[
 \mu_r\leq Z/(2r)+1\leq2Z/r
\]
integers in one residue class modulo $r$.  With $M=|B|\gg Z/\log Z$, \cref{lem:cyclic-energy} gives
\[
 D_r\gg\frac{M^3}{\mu_r^2r^2}
 \gg\frac{Z^3/\log^3Z}{Z^2}
 =\frac{Z}{\log^3Z}.
\]
Moreover $M/\mu_r\gg r/\log Z\geq X/\log Z$, so the largeness hypothesis is uniform in $r$.
\end{proof}

\begin{proposition}[Target-row distance]
\label{prop:target-row-distance}
For every $r\in S_b$,
\[
 D_r\geq\frac{|B|}{r^2}\gg_b\frac{Z}{\log Z}.
\]
\end{proposition}

\begin{proof}
Every $q\in B$ is coprime to $r$.  If $d\not\equiv0\pmod r$, then $dq^{-1}$ is nonzero modulo $r$, so each circle distance is at least $1/r$.
\end{proof}

The target-row sensors depend only on the reduced denominator $b$.  Their common kernel on $\prod_{r\in S_b}\mathbb Z/r\mathbb Z$ is trivial.  The numerator $a$ of the fixed target does not enter these sensors and creates no additional decoder or CRT obstruction.

\subsection{Row tails}

For $r\in R_{\mathrm{rows}}$, define
\[
 E_r(x)=\sum_{q\in B}\dist{u_{rq}(x)}^2,
 \qquad
 A_r(x)=\prod_{q\in B}
 \left|(1-\theta)+\theta\exp(2\pi i u_{rq}(x))\right|,
\]
and choose $x_r^*$ minimizing $E_r(x)$.

\begin{lemma}[Shift-uniform row tail]
\label{lem:row-tail}
There is
\[
 c_{\rm row}=c_{\rm row}(t,\gamma,b)>0
\]
such that
\[
 \delta_r:=\sum_{x\neq x_r^*}A_r(x)
 \leq r\exp(-c_{\rm row}D_r).
\]
\end{lemma}

\begin{proof}
If $x\neq x_r^*$, \cref{eq:row-shift} and the circle triangle inequality imply
\[
 D_r\leq2E_r(x)+2E_r(x_r^*)\leq4E_r(x).
\]
Apply \cref{eq:bernoulli-modulus} with $\kappa_{\rm mod}$ and sum over the at most $r$ nondecoder residues.  One may take $c_{\rm row}=\kappa_{\rm mod}/4$.
\end{proof}

Consequently,
\begin{equation}
\label{eq:Delta-row}
 \Delta:=\sum_{r\in R_{\mathrm{rows}}}\delta_r
 \leq Z^2\exp(-c_{\rm row}Z/\log^3Z)=o(1).
\end{equation}

\subsection{The exact retained skeleton}

For the fixed anchor assignment, assign every factor involving one anchor prime and one row prime to the corresponding row modulus $A_r$.  Internal anchor factors occur in $T_B$, while every retained lower--lower factor and every phase remains in a residual factor.  Thus
\begin{equation}
\label{eq:concrete-fibre-factorization}
 F(a_B,x_R)
 =T_B(a_B)\prod_{r\in R_{\mathrm{rows}}}A_r(x_r)
  H_{a_B}(x_R),
 \qquad |H_{a_B}(x_R)|\leq1.
\end{equation}
If the full modulus vanishes, set $H_{a_B}=0$.  The exact weighted compression theorem and \cref{eq:P-top,eq:Delta-row} give:

\begin{proposition}[Global weighted fibre error]
\label{prop:global-fibre-error}
The total contribution of every nondecoder fibre point is at most
\begin{equation}
\label{eq:global-fibre-error}
 P_{\mathrm{anc}}\bigl(\exp(\Delta)-1\bigr)
 \ll Z\log Z\,\Delta=o(1).
\end{equation}
\end{proposition}

No decoder weight is divided out, and the raw number of anchor assignments does not occur.

\subsection{Identification ranges}

Suppose the anchor assignment is coherent with integer label $m$.  The candidate row residue $m\pmod r$ satisfies
\begin{equation}
\label{eq:candidate-row-energy}
 E_r(m\bmod r)
 \leq\frac{m^2}{r^2}W_B
 \ll\frac{m^2/r^2}{Z\log Z}.
\end{equation}
If this is less than $D_r/8$, the candidate is the unique minimizer.

Put
\begin{equation}
\label{eq:M-dec}
 M_{\mathrm{dec}}=\frac{XZ}{(\log Z)^2}.
\end{equation}

\begin{proposition}[Prime-coordinate identification]
\label{prop:prime-decoder-identification}
For every $r\in P\setminus B$ and every coherent anchor label satisfying $|m|\leq M_{\mathrm{dec}}$,
\[
 x_r^*=m\pmod r.
\]
\end{proposition}

\begin{proof}
At $r=X$ and $|m|=M_{\mathrm{dec}}$, \cref{eq:candidate-row-energy} is $O(Z/\log^5Z)$, whereas \cref{prop:lower-row-distance} gives $D_r\gg Z/\log^3Z$.
\end{proof}

Fix once and for all a sufficiently small constant
\[
 \kappa_0=\kappa_0(t,\gamma,b)>0,
 \qquad \kappa_0<\frac14,
\]
small enough for the finitely many target rows and for the common linear phase range below.  Define the universal full-coordinate cutoff
\begin{equation}
\label{eq:T-full}
 T_0=\kappa_0\min(X^2,Z).
\end{equation}

\begin{proposition}[Target-coordinate identification]
\label{prop:target-decoder-identification}
For every $r\in S_b$ and every coherent anchor label satisfying $|m|\leq T_0$,
\[
 x_r^*=m\pmod r.
\]
\end{proposition}

\begin{proof}
By \cref{eq:candidate-row-energy}, the candidate energy is $O_b(m^2/(Z\log Z))$, while \cref{prop:target-row-distance} is $\gg_b Z/\log Z$.  Since $|m|\leq T_0\leq\kappa_0Z$, choosing $\kappa_0$ sufficiently small makes the comparison strict, uniformly over the finite set $S_b$.
\end{proof}

Because $b$ is squarefree, agreement modulo every $r\in S_b$ is agreement modulo $b$.  Hence for $|m|\leq T_0$, the entire decoded tuple is the genuine integer $m$ modulo $L$.  In the larger range $T_0<|m|\leq M_{\mathrm{dec}}$, all $P$-coordinates still decode to $m$, while no target-coordinate identification is asserted.  The transition and adaptive estimates use only those prime coordinates.

\section{The decoded skeleton and its minor ranges}
\label{sec:decoded-skeleton}

After \cref{prop:global-fibre-error}, it remains to estimate the single decoded tuple attached to each anchor assignment.  Coherent anchor assignments are indexed by an integer label $m$; noncoherent assignments are treated by the anchor partition.  The lower--lower pair factors were retained during fibre compression.

\subsection{Prime-only Gaussian transition}

Assume $1<\gamma\leq2$ and
\begin{equation}
\label{eq:transition-range}
 T_{\mathrm{full}}<|m|\leq X^2/4.
\end{equation}
By \cref{prop:prime-decoder-identification}, every coordinate indexed by $P$ equals $m$.  Hence every complete-pair factor has phase $m/(pq)$.  Since $pq\geq X^2$ and $|m|\leq X^2/4$, these phases lie in the linear circle range, and \cref{eq:bernoulli-modulus} gives
\begin{equation}
\label{eq:transition-pointwise}
 |F(h_m)|
 \leq\exp\left(-c m^2
   \sum_{p<q\in P}\frac1{p^2q^2}\right)
 \leq\exp\left(-c\frac{m^2}{X^2\log^2X}\right).
\end{equation}
At the lower endpoint, the exponent satisfies
\begin{equation}
\label{eq:transition-lower-exponent}
 \frac{T_{\mathrm{full}}^2}{X^2\log^2X}
 \gg\frac{Z^2}{X^2\log^2X}
 =\frac{X^{2\gamma-2}}{\log^2X}\longrightarrow\infty.
\end{equation}
Therefore
\begin{equation}
\label{eq:transition-total}
 \sum_{T_{\mathrm{full}}<|m|\leq X^2/4}|F(h_m)|
 =o(\sigma_{\mathcal E}^{-1}).
\end{equation}
This sector uses no target coordinate.  At $\gamma=2$ it treats the variance boundary exactly; for $1<\gamma<2$ it fills the entire gap left by the shorter full-coordinate range.

\subsection{Adaptive retained-pair damping}

For an integer $m$ with
\[
 X^2/4<|m|\leq M_{\mathrm{dec}},
\]
put $I_m=[2\sqrt{|m|},3\sqrt{|m|}]$.

\begin{proposition}[Adaptive prime interval]
\label{prop:adaptive-prime-interval}
For every such $m$, the interval $I_m$ lies in $[X,Z/2)$ and contains
\[
 K_m\gg\frac{\sqrt{|m|}}{\log|m|}
\]
primes.  On the decoded skeleton, the retained complete-pair energy satisfies
\begin{equation}
\label{eq:adaptive-energy}
 Q_{\mathrm{pair}}(m)\gg\frac{|m|}{(\log|m|)^2}.
\end{equation}
\end{proposition}

\begin{proof}
The lower inclusion follows from $2\sqrt{|m|}>X$.  At the upper endpoint,
\[
 3\sqrt{|m|}\leq\frac{3\sqrt{XZ}}{\log Z}<Z/2
\]
for every fixed $\gamma>1$ and large $X$.  Fixed-ratio prime supply gives the count.  For distinct $p,q\in I_m$, one has $1/9\leq|m|/(pq)\leq1/4$, so each retained pair contributes a fixed positive amount of energy; summing over the pairs proves \cref{eq:adaptive-energy}.
\end{proof}

Consequently,
\begin{equation}
\label{eq:adaptive-pointwise}
 |F(h_m)|\leq\exp\left(-c\frac{|m|}{(\log|m|)^2}\right),
\end{equation}
and
\begin{equation}
\label{eq:adaptive-total}
 \sum_{X^2/4<|m|\leq M_{\mathrm{dec}}}|F(h_m)|
 =o(\sigma_{\mathcal E}^{-1}).
\end{equation}

\subsection{Pairwise-disjoint exhaustive partitions}

Fix $C\geq1$, to be chosen before $X$, and put
\begin{equation}
\label{eq:N-major}
 N=\left\lfloor\frac{C}{\sigma_{\mathcal E}}\right\rfloor.
\end{equation}
By \cref{prop:variance-regimes}, for every fixed $C$,
\begin{equation}
\label{eq:N-below-full}
 N=o(T_{\mathrm{full}}).
\end{equation}

Apply the decoder/nondecoder split first.  Every nondecoder tuple belongs to the terminal error sector.  Among decoded tuples, a noncoherent anchor assignment belongs to that same sector.  A coherent assignment has one unique label $m$.

If $\gamma>2$, the following five sectors are pairwise disjoint and exhaustive:
\begin{center}
\begin{tabular}{@{}c p{5.2cm} p{5.0cm}@{}}
\toprule
Sector & skeleton condition & control \\
\midrule
I & $|m|\leq N$ & positive Taylor major \\
II & $N<|m|\leq X^2/4$ & actual-variance Gaussian tail \\
III & $X^2/4<|m|\leq M_{\mathrm{dec}}$ & adaptive retained pairs \\
IV & $|m|>M_{\mathrm{dec}}$ & internal anchor energy \\
V & noncoherent decoded tuple, or any nondecoder tuple & anchor entropy and fibre error \\
\bottomrule
\end{tabular}
\end{center}

If $1<\gamma\leq2$, the following six sectors are pairwise disjoint and exhaustive:
\begin{center}
\begin{tabular}{@{}c p{5.2cm} p{5.0cm}@{}}
\toprule
Sector & skeleton condition & control \\
\midrule
I & $|m|\leq N$ & positive Taylor major \\
II & $N<|m|\leq T_{\mathrm{full}}$ & actual-variance Gaussian tail \\
III & $T_{\mathrm{full}}<|m|\leq X^2/4$ & prime-only pair Gaussian \\
IV & $X^2/4<|m|\leq M_{\mathrm{dec}}$ & adaptive retained pairs \\
V & $|m|>M_{\mathrm{dec}}$ & internal anchor energy \\
VI & noncoherent decoded tuple, or any nondecoder tuple & anchor entropy and fibre error \\
\bottomrule
\end{tabular}
\end{center}

\begin{proposition}[Regime-dependent sector exhaustion]
\label{prop:five-sector-exhaustion}
The five sectors for $\gamma>2$, and the six sectors for $1<\gamma\leq2$, are respectively pairwise disjoint and exhaustive over all frequencies modulo $L$.
\end{proposition}

\begin{proof}
The CRT coordinate map is a bijection.  The decoder/nondecoder split is disjoint for each anchor assignment.  A decoded coherent assignment has exactly one label and therefore lies in exactly one of the displayed disjoint label ranges.  No residual sibling sector remains.
\end{proof}

\subsection{Outer totals}

For $|m|>M_{\mathrm{dec}}$, the internal anchor factor gives
\[
 |F(h_m)|\leq\exp(-\kappa_bm^2\sigma_{B,0}^2).
\]
Thus
\begin{equation}
\label{eq:anchor-tail-total}
 \sum_{|m|>M_{\mathrm{dec}}}|F(h_m)|
 \ll\sigma_{B,0}^{-1}
 \exp(-cM_{\mathrm{dec}}^2\sigma_{B,0}^2)
 =o(\sigma_{\mathcal E}^{-1}),
\end{equation}
because $M_{\mathrm{dec}}^2\sigma_{B,0}^2\gg X^2/\log^6Z$.

Finally, \cref{eq:noncoherent-top-tail,eq:global-fibre-error} give
\begin{equation}
\label{eq:terminal-error-total}
 \sum_{\mathrm{terminal}}|F|
 \ll\exp(-cF_B)+P_{\mathrm{anc}}(\exp(\Delta)-1)
 =o(\sigma_{\mathcal E}^{-1}).
\end{equation}

\section{The major contribution and terminal budget}
\label{sec:major-budget}

We now estimate the Taylor major and the full-variance tail, then assemble the regime-dependent sector partition.  The order of choices is essential: the Gaussian cutoff $C$ is fixed first, and only then is $X$ taken sufficiently large.  All constants are fixed once $t$, $\gamma$, $b$ and the stated numerical choices are fixed.

\subsection{A uniform local expansion}

Compactness of the Bernoulli interval $I_{t,\gamma}$ gives constants $\rho>0$ and $M<\infty$ such that, uniformly for $\theta\in I_{t,\gamma}$ and $|z|\leq\rho$,
\begin{equation}
\label{eq:local-log-expansion}
 \log\bigl((1-\theta)+\theta\exp(2\pi iz)\bigr)
 =2\pi i\theta z
 -2\pi^2\theta(1-\theta)z^2
 +R_\theta(z),
\end{equation}
where
\begin{equation}
\label{eq:local-log-remainder}
 |R_\theta(z)|\leq M|z|^3.
\end{equation}

Let
\[
 e_{\min}=\min_{e\in\mathcal E}e\asymp\min(X^2,Z).
\]
For $|m|\leq N=C/\sigma_{\mathcal E}$,
\begin{equation}
\label{eq:aggregate-cubic-general}
 \sum_{e\in\mathcal E}|m/e|^3
 \leq\frac{|m|^3}{e_{\min}}
       \sum_{e\in\mathcal E}\frac1{e^2}
 \ll\frac{C^3}{\sigma_{\mathcal E}e_{\min}}.
\end{equation}
By \cref{prop:variance-regimes}, this is
\begin{equation}
\label{eq:aggregate-cubic-remainder}
 \begin{cases}
 O(C^3\log X/X),&\gamma>2,\\
 O(C^3\sqrt{\log Z/Z}),&1<\gamma\leq2,
 \end{cases}
\end{equation}
and hence tends to zero for fixed $C$ in every regime.  Also $N=o(e_{\min})$, so all major phases lie in the common Taylor disk.

\subsection{Sector I: positive real major}

For $|m|\leq N$, \cref{eq:N-below-full} and the decoder-identification propositions show that the decoded tuple is the genuine integer frequency $m$ modulo $L$.  Summing \cref{eq:local-log-expansion} over the actual family, the linear term cancels exactly:
\[
 2\pi i m\theta\Lambda-2\pi i mt=0.
\]
The numerator $a$ has therefore introduced no residual phase.  The quadratic term is exactly the full variance.  Uniformly for $|m|\leq N$,
\begin{equation}
\label{eq:major-local-gaussian}
 F(m)=\exp\bigl(-2\pi^2m^2\sigma_{\mathcal E}^2+\varepsilon_m\bigr),
 \qquad
 \max_{|m|\leq N}|\varepsilon_m|\longrightarrow0.
\end{equation}
Thus, for all sufficiently large $X$,
\[
 \RePart F(m)\geq c_0\exp(-2\pi^2m^2\sigma_{\mathcal E}^2)
\]
with $c_0>0$ fixed.  Summing on $|m|\leq1/(4\sigma_{\mathcal E})$ gives:

\begin{proposition}[Positive Gaussian major]
\label{prop:positive-major}
For every fixed $C\geq1$ and all sufficiently large $X$,
\begin{equation}
\label{eq:positive-major}
 \RePart\sum_{|m|\leq N}F(m)
 \geq\frac{c_{\mathrm{maj}}}{\sigma_{\mathcal E}},
\end{equation}
where $c_{\mathrm{maj}}>0$ is independent of $C$ and $X$.
\end{proposition}

\subsection{Sector II: the full actual-variance tail}

For
\[
 N<|m|\leq T_{\mathrm{full}},
\]
the decoded tuple is the genuine integer $m$ modulo $L$.  The constant in \cref{eq:T-full} was chosen so that $|m|/e\leq1/4$ for every $e\in\mathcal E$.  Hence
\begin{equation}
\label{eq:full-variance-gaussian}
 |F(m)|\leq\exp(-c m^2\sigma_{\mathcal E}^2).
\end{equation}
A Gaussian integral comparison gives
\begin{equation}
\label{eq:sector-II-tail}
 \sum_{|m|>N}\exp(-c m^2\sigma_{\mathcal E}^2)
 \leq K\sigma_{\mathcal E}^{-1}\exp(-cC^2/2),
\end{equation}
where $K$ is independent of $C$ and $X$.  Choose $C$ so large that this is at most $c_{\mathrm{maj}}/(8\sigma_{\mathcal E})$.

\subsection{Strict assembly in all regimes}

Fix this $C$, then choose $X$ above the finite maximum of the thresholds required for:
\begin{itemize}
\item the parameterized prime counts, reciprocal masses, family size and complete square load;
\item separation from the support of $b$ and avoidance of $\mathcal T$;
\item the strict load inequalities $t<\Lambda<1$ and the compact Bernoulli interval;
\item anchor rigidity, fingerprint entropy and the weighted anchor partition;
\item both row-distance estimates, the row-tail sum and the two identification ranges;
\item the prime-only transition estimate when $\gamma\leq2$;
\item containment and prime supply in the adaptive interval;
\item the Taylor disk and the regime-specific cubic remainder;
\item all $o(\sigma_{\mathcal E}^{-1})$ totals in \cref{eq:transition-total,eq:adaptive-total,eq:anchor-tail-total,eq:terminal-error-total}.
\end{itemize}
The terminal order is therefore
\begin{equation}
\label{eq:parameter-order}
 (t,\gamma,b,\mathcal T)\ \text{fixed},
 \qquad C\longrightarrow X.
\end{equation}
For $\gamma>2$, Sectors III--V contribute $o(\sigma_{\mathcal E}^{-1})$.  For $1<\gamma\leq2$, Sectors III--VI do so, including the prime-only transition.  Sector II is at most $c_{\mathrm{maj}}/(8\sigma_{\mathcal E})$, whereas Sector I has real part at least $c_{\mathrm{maj}}/\sigma_{\mathcal E}$.  Applying \cref{prop:decoded-skeleton-positivity} yields:

\begin{theorem}[Positive parameterized Fourier numerator]
\label{thm:positive-fourier-numerator}
For every fixed reduced $t=a/b\in(0,1)$ with squarefree $b$, every fixed $\gamma$ satisfying \cref{eq:admissible-region}, and every finite forbidden set $\mathcal T$, one has for all sufficiently large $X$
\begin{equation}
\label{eq:positive-full-sum}
 \RePart\sum_{h\bmod L}F(h)>0.
\end{equation}
\end{theorem}

The proof uses the full actual-family variance in both the major and the first minor range.  At $\gamma=2$, that variance is target-row dominated, while the following transition range is controlled solely by the complete prime-pair square load.  No estimate is obtained by formal substitution for the exponent $3$.

\section{Exact representation and completion of the theorem}
\label{sec:exact-completion}

We now pass from the positive finite Fourier coefficient to an exact rational identity, and then complete the reductions announced in \cref{sec:introduction}.

\subsection{The avoiding unit representation}

Let $W$ be the Bernoulli weight of subsets $A\subset\mathcal E$ satisfying
\[
 \sum_{e\in A}\frac{L}{e}\equiv\frac{L}{b}\pmod L.
\]
By \cref{eq:reciprocal-fourier,thm:positive-fourier-numerator},
\[
 LW=\sum_{h\bmod L}F(h)
\]
has positive real part.  The left side is real and nonnegative, so $W>0$.  Hence there exists a subset $A\subset\mathcal E$ with
\begin{equation}
\label{eq:target-congruence}
 \sum_{e\in A}\frac{L}{e}\equiv\frac{L}{b}\pmod L.
\end{equation}
Dividing by $L$, the difference
\[
 d_A=\sum_{e\in A}\frac1e-\frac1b
\]
is an integer.

The exactness step is now deterministic.  By \cref{eq:exact-centering-load},
\[
 0\leq\sum_{e\in A}\frac1e\leq\Lambda<1,
 \qquad
 0<1/b<1.
\]
Therefore
\[
 -1<d_A<1.
\]
Since $d_A\in\mathbb Z$, it follows that $d_A=0$.  Thus
\begin{equation}
\label{eq:avoiding-unit-identity}
 \sum_{e\in A}\frac1e=\frac1b.
\end{equation}
By \cref{prop:arithmetic-capacity}, the elements of $A$ are distinct squarefree semiprimes and avoid the prescribed finite set $\mathcal T$.

We have proved the fixed-denominator statement on which the elementary closure rests.

\begin{theorem}[Avoiding unit representation]
\label{thm:avoiding-unit}
Let $b\geq3$ be squarefree and let $\mathcal T$ be a finite set of positive integers.  There is a finite set $A$ of pairwise distinct squarefree semiprimes, disjoint from $\mathcal T$, such that
\[
 \sum_{e\in A}\frac1e=\frac1b.
\]
\end{theorem}

This is the unique point at which quotient positivity and ambient equality meet.  The Fourier argument proves the congruence \cref{eq:target-congruence}; only afterwards does the strict interval $[0,\Lambda]\subset[0,1)$ exclude every nonzero alias.

\subsection{Numerators and small denominators}

\begin{lemma}[Numerator induction]
\label{lem:numerator-induction}
Suppose an avoiding representation of $1/b$ exists for every finite forbidden set.  Then, for every integer $a\geq1$ and every finite forbidden set $\mathcal T$, there is an avoiding representation of $a/b$ by distinct squarefree semiprimes.
\end{lemma}

\begin{proof}
Induct on $a$.  The case $a=1$ is the hypothesis.  Suppose a finite set $A$ represents $a/b$ and avoids $\mathcal T$.  Apply the avoiding unit statement with forbidden set $\mathcal T\cup A$ to obtain a representation $A'$ of $1/b$ disjoint from $A$.  Then $A\cup A'$ consists of distinct squarefree semiprimes, avoids $\mathcal T$, and has reciprocal sum $(a+1)/b$.
\end{proof}

\begin{lemma}[Denominators $1$ and $2$]
\label{lem:small-denominators}
It is enough to prove \cref{thm:avoiding-unit} for squarefree $b\geq3$.
\end{lemma}

\begin{proof}
For $b=2$, construct successively disjoint avoiding representations of $1/3$ and $1/6$, and use
\[
 \frac12=\frac13+\frac16.
\]
For $b=1$, construct three successive avoiding representations of $1/3$ and add them; equivalently, apply \cref{lem:numerator-induction} with numerator and denominator equal to $3$.  At each stage the previously used denominators are added to the forbidden set.  Once an avoiding unit representation for $1$ or $1/2$ is obtained, \cref{lem:numerator-induction} supplies every positive numerator.
\end{proof}

\subsection{Necessity of squarefreeness}

\begin{lemma}[Squarefree denominator obstruction]
\label{lem:squarefree-necessity}
Let $A$ be a finite set of squarefree positive integers.  Then the reduced denominator of
\[
 \sum_{n\in A}\frac1n
\]
is squarefree.
\end{lemma}

\begin{proof}
Let $M=\operatorname{lcm}(A)$.  Then
\[
 \sum_{n\in A}\frac1n
 =\frac{\sum_{n\in A}M/n}{M}.
\]
The reduced denominator divides $M$.  Since the least common multiple of squarefree integers is squarefree, every divisor of $M$ is squarefree.
\end{proof}

\begin{proof}[Proof of \cref{thm:main}]
If $q$ is a sum of reciprocals of distinct squarefree semiprimes, its reduced denominator is squarefree by \cref{lem:squarefree-necessity}.

Conversely, write $q=a/b>0$ in lowest terms and suppose $b$ is squarefree.  For $b\geq3$, \cref{thm:avoiding-unit,lem:numerator-induction} give a representation of $a/b$ by distinct squarefree semiprimes.  The cases $b=1,2$ follow from \cref{lem:small-denominators}.  This proves the characterization.
\end{proof}

\subsection{Analytic input}

The prime number theorem is the only external analytic result used in the construction.  It supplies prime counts in fixed-ratio intervals and, by Abel summation, the reciprocal load of the prime blocks.  The synchronization, fibre compression, target observability, five-sector decomposition, positive major estimate, no-wrap step and elementary induction are all proved within the article.

\section{Structural scope and limits}
\label{sec:scope}

The proof was developed for a specific reciprocal-representation problem, but several of its mechanisms are independent of primes and semiprimes.  This section records their natural scope without turning the article into a separate abstract-framework paper.

\subsection{What is genuinely reusable}

The exact weighted product-fibre theorem, \cref{thm:weighted-compression}, is a deterministic statement about finite product spaces.  It requires no CRT structure, no probability interpretation and no arithmetic input.  Its main content is the normalization: off-decoder mass is averaged against the actual anchor weight, not multiplied by the number of anchor assignments.  The tensorized expression is optimal from the displayed row data alone.  It also preserves every unassigned factor on the decoded skeleton, where that factor may continue to supply cancellation or modulus decay.

The syndrome theorem, \cref{thm:syndrome-separation,prop:qualitative-observability}, gives a finite-group language for hidden-coordinate sensitivity.  Its invariant is the pseudometric
\[
 d_\Psi(h,h')^2
 =\sum_j\dist{\psi_j(h-h')}^2.
\]
The zero-distance subgroup is precisely the common kernel of the sensors.  In the present application the sensors are the top-cross rows, and squarefreeness permits the target group to decompose into prime coordinates, each of which is detected directly.

The decoded-skeleton positivity proposition, \cref{prop:decoded-skeleton-positivity}, is an exact assembly interface.  It allows any finite collection of minor lanes and does not impose a universal distinction between major arcs, middle arcs, energetic anchors or outer frequencies.  Its proof is deliberately short: the difficult work lies in producing the anchor partition, decoders and lane estimates in a concrete problem.

Finally, \cref{thm:complete-family-collision} separates complete-family mass from rigidity.  It applies to every nonnegative weighted complete $k$-uniform family.  In the present proof only the exact pair identity is needed, but the collision inequality identifies the moving-uniformity range
\[
 k^2V/U^2\longrightarrow0
\]
in which the first two power sums control the complete-family mass.

\subsection{Product measures beyond Bernoulli selectors}

The Fourier product formula in \cref{lem:independent-product} extends immediately to independent finite-valued random variables.  Each local Bernoulli factor is replaced by the Fourier transform of the corresponding local probability law.  Whenever the character space admits an anchor--fibre factorization and the local transforms satisfy suitable decoder and residual estimates, the same weighted compression and skeleton assembly apply.

The algebraic decomposition also makes sense for signed or complex local weights after absolute values are assigned to the row kernels and phases are retained in the residual factor.  The probabilistic interpretation must then be replaced by a statement about a nonzero or positive coefficient.  We do not pursue this extension here, because it has no additional role in the semiprime application.

\subsection{Averaged rather than worst-case fibres}

The normalized corollary \cref{cor:normalized-fibre} permits the exact anchor-dependent error
\[
 \sum_yA_y\left\{\prod_r(1+\varepsilon_{r,y})-1\right\}.
\]
Thus a small set of poorly decoded anchors is harmless when its decoded anchor mass is small.  The E306 proof uses a uniform row-tail estimate and then the global top partition, but the general statement does not require one worst-case value of $\Delta_y$.

This flexibility is useful only when the product structure is preserved.  A genuinely sequential kernel
\[
 a_r(x_r\mid x_1,\ldots,x_{r-1},y)
\]
would require a different conditional-product theorem; no such theorem is asserted here.

\subsection{Failure modes}

Three failures are structural rather than quantitative.

\paragraph{Unweighted fibre counting.}
Multiplying a row error by the raw cardinality of the anchor space may destroy the estimate or simply be false.  The correct object is the weighted tensorized error in \cref{eq:exact-fibre-error}.

\paragraph{Discarding residual factors.}
Bounding every unassigned Fourier factor by $1$ before reaching the decoded skeleton can remove the only available damping in an intermediate label range.  Sector III is possible precisely because the lower--lower pair factors remain visible after row decoding.

\paragraph{Unsensed target coordinates.}
If a nontrivial quotient direction lies in the common kernel of all sensors and residual factors, its nontrivial character may cancel an apparently positive visible contribution.  Target sensitivity is therefore a coverage obligation: in each lane, every hidden direction must be fixed, decoded or residually damped.

\subsection{Present limitations}

The compression theorem takes absolute values on off-decoder fibres.  It does not exploit cancellation among several nondecoder tuples or among different fibres.  A theorem using genuinely complex fibre cancellation would be a mechanism-level strengthening rather than a sharper provider estimate.

Likewise, complete-family mass does not imply higher-uniformity synchronization.  For pairs, the graph of denominator factors interacts naturally with one-anchor row codes.  A $k$-uniform extension would require new rigidity and decoding results; the collision estimate alone does not provide them.

At present, Erd\H{o}s Problem 306 is the first substantive consumer of the full anchor--fibre organization.  A separate framework paper would become mathematically justified only after a second independent application or a genuine strengthening such as complex fibre cancellation, non-enumerative synchronization, or a broader observability classification.  Until then the general tools belong naturally inside the proof which produced and uses them.

\subsection{Relation to formal verification}

A separate Lean development establishes the headline representation theorem along an orthogonal formalization axis.  Its analytic trust boundary and proof decomposition are not identical to the dense one-anchor human argument presented here.  Accordingly, the formal development is evidence for the theorem but is not invoked to certify any step of this ordinary proof.  The present article exposes its own complete dependency chain and may be checked independently of the formal source.

\bibliographystyle{alpha}
\bibliography{references}

\end{document}
